\chapter*{前言}

很多人学 \cpp 的痛苦不在语法本身，而在语法背后的执行模型没有连起来：变量到底放在哪里，引用为什么会悬空，构造函数什么时候运行，模板报错为什么像墙一样长，标准库容器为什么有时快、有时慢。只背语法表，很快就会写出能编译但不可维护的程序；只看项目，又会被概念淹没。

这本书按另一条路写：每章先从一个具体小问题开始，再把代码、内存、类型、接口和工程约束一起讲清楚。它不是语法速查表，也不是只讲高级技巧的炫技书。它的目标是让读者从第一段程序走到能写可测试、可维护、可解释性能的现代 \cpp 代码。

书中代码默认使用 \cpp20。示例会优先使用标准库、RAII、值语义、明确所有权和小接口；类成员函数中的成员访问会写出 \code{this->}，这是为了在教学阶段把“成员”和“局部变量”区分清楚。运行时代码尽量避免递归，树和图的处理会用显式栈、队列或工作表说明，因为工程代码里递归深度常常不是一个可以随手忽略的细节。

\begin{keyidea}
学 \cpp 不是记住更多语法，而是建立三条主线：第一，值和对象如何存在；第二，资源如何被拥有和释放；第三，接口如何把复杂性限制在可测试的边界内。
\end{keyidea}
